% Part: computability
% Chapter: computability-theory
% Section: complete-ce-sets

\documentclass[../../../include/open-logic-section]{subfiles}

\begin{document}

\olfileid{cmp}{thy}{cce}
\olsection{संपूर्ण संगणनक्षमपणे प्रगणनीय संच}

\begin{defn}
संच $A$ हा (अनेक-एक न्यूनीकरणक्षमतेखाली) \emph{संपूर्ण
!!{computably enumerable} संच} असतो, जर
\begin{enumerate}
\item $A$ संगणनक्षमपणे प्रगणनीय असेल, आणि
\item इतर कोणताही संगणनक्षमपणे प्रगणनीय संच $B$ साठी $B \leq_m A$ असेल.
\end{enumerate}
\end{defn}

म्हणजे, संपूर्ण संगणनक्षमपणे प्रगणनीय संच हे अशा संचांमधील
शक्य तितके ``कठीण'' संच असतात. त्यांच्याद्वारे
\emph{कोणत्याही} संगणनक्षमपणे प्रगणनीय संचाविषयीच्या
प्रश्नांची उत्तरे मिळवता येतात.

\begin{thm}
$K$, $K_0$, आणि $K_1$ हे तिन्ही संपूर्ण संगणनक्षमपणे
प्रगणनीय संच आहेत.
\end{thm}

\begin{proof}
$K_0$ संपूर्ण आहे हे पाहण्यासाठी $B$ हा कोणताही
संगणनक्षमपणे प्रगणनीय संच असू द्या. मग काही निर्देशांक $e$
साठी
\[
B = W_e = \Setabs{x}{\cfind{e}(x) \fdefined}.
\]
$f(x) = \tuple{e, x}$ हे फलन $f$ असू द्या. मग प्रत्येक
नैसर्गिक संख्या $x$ साठी $x \in B$ असेल तेव्हाच
$f(x) \in K_0$. म्हणजे $f$ हे $B$ चे~$K_0$ कडे
न्यूनीकरण करते.

$K_1$ संपूर्ण आहे हे पाहण्यासाठी, \olref[k1]{prop:k1}
च्या सिद्धतेत आपण $K_0$ चे त्याच्याकडे न्यूनीकरण केले
आहे, हे लक्षात घ्या. म्हणून \olref[ppr]{prop:trans-red}
नुसार कोणताही संगणनक्षमपणे प्रगणनीय संच~$K_1$ कडेही
न्यूनीत करता येतो.

$K_0$ चे~$K$ कडेही अशाच प्रकारे न्यूनीकरण करता येते.
%READERNOTE{OLCMP-023 — गोठवलेल्या इंग्रजीतील शेवटचे सिद्धतावाक्य K चे K_0 कडे न्यूनीकरण सांगते. त्यामुळे K_0 च्या संपूर्णतेवरून K ची संपूर्णता सिद्ध होत नाही; आवश्यक दिशा K_0 चे K कडे अशी आहे. पुढील स्वतंत्र प्रश्न मात्र मूळ K चे K_0 कडे न्यूनीकरणच विचारतो, आणि तो तसाच ठेवला आहे.}
\end{proof}

\begin{prob}
$K$ चे $K_0$ कडे न्यूनीकरण द्या.
\end{prob}

\begin{digress}
आतापर्यंत आपण पाहिलेल्या संगणनक्षमपणे प्रगणनीय संचांची सर्व
उदाहरणे एकतर संगणनक्षम किंवा संपूर्ण आहेत. हे विचित्र वाटले
पाहिजे!{} संगणनक्षमही नाहीत आणि संपूर्णही नाहीत, अशी
संगणनक्षमपणे प्रगणनीय संचांची उदाहरणे आहेत का? उत्तर होय
आहे; मात्र 1950 च्या दशकाच्या मध्यापर्यंत हे सिद्ध झाले
नव्हते. फ्रीडबर्ग आणि मुछ्निक यांनी स्वतंत्रपणे ते सिद्ध केले.
\end{digress}

\end{document}
